\section*{Resumen}
La colocación continua de una fundación de hormigón masivo acorta la jornada de construcción, pero obliga a controlar el calor que se acumula en el interior y las diferencias de temperatura que aparecen al enfriarse sus superficies. Se estudió una fundación de aerogenerador de $636{,}775\,\mathrm{m^3}$ para el sitio docente de Los Ángeles, Chile. Un modelo de volúmenes finitos representó la geometría del proyecto, la colocación en once capas, las protecciones y nueve circuitos de refrigeración. Se combinaron ocho curvas de hidratación obtenidas de ensayos de probeta con tres historias meteorológicas del reanálisis ERA5. Entre los planes evaluados, la colocación continua a $60\,\mathrm{m^3/h}$ terminó en $10{,}61\,\mathrm h$, unos $58$ minutos antes que la alternativa de $55\,\mathrm{m^3/h}$. Con la última manta hasta los $21$ días de edad de cada capa y sombra posterior, ninguno de los $24$ escenarios superó los límites docentes de $70\,^\circ\mathrm C$ o $20\,\mathrm K$ durante los $28$ días de seguimiento. Con sombra y retiro a los $15$ días, cinco escenarios superaron al menos un límite en la malla de comparación. La recomendación depende de que en obra se alcancen la temperatura de colocación, la refrigeración y la protección solar previstas.

\textbf{Palabras clave}\quad hormigón masivo, calor de hidratación, control térmico, volúmenes finitos, planificación de la colocación
